\documentclass[10pt,a4paper]{amsart}
\usepackage[margin=19mm]{geometry}
\usepackage{amsmath,amssymb,mathtools}
\usepackage[scr=rsfs,cal=euler]{mathalfa}
\usepackage{xfrac}
\usepackage[unicode,hidelinks,bookmarks=false]{hyperref}
\newcommand{\ie}{i.e.\ }
\newcommand{\cf}{cf.\ }
\newcommand{\resp}{respectively\ }
\newcommand{\Subsection}{\subsection}
\providecommand{\nobreakdash}{\nobreak\hbox{-}}
\setlength{\emergencystretch}{2em}
\multlinegap0pt
\usepackage{bm}
\usepackage{bbm}
\usepackage{dsfont}
\usepackage{xspace}
\usepackage{cancel}
\usepackage{mathtools}
\usepackage{amsthm}
\makeatletter
\@ifundefined{Definition}{\newtheorem{Definition}{Definition}}{}
\@ifundefined{Lemma}{\newtheorem{Lemma}[Definition]{Lemma}}{}
\@ifundefined{Proposition}{\newtheorem{Proposition}[Definition]{Proposition}}{}
\makeatother
\makeatletter
\newcommand{\R}{\mathbb R}
\expandafter\ifx\csname inlineenum\endcsname\relax
\newenvironment{inlineenum}
 {\setcounter{inlineenum}{0}%
  \renewcommand{\item}{\refstepcounter{inlineenum}{(\theinlineenum)\hspace{\thelabelsep}}}
 }
\fi
\makeatother
\title[Failure of local equi-Lipschitzness for families of Lorentz ]{Failure of local equi-Lipschitzness for families of Lorentz distances to Cauchy surface foliations}
\author{Gregory J. Galloway, Robert J. McCann, Argam Ohanyan}
\date{Source: arXiv:2606.21997v2 (CC BY 4.0)}
\begin{document}
\maketitle

\noindent\emph{Source and licence.} Failure of local equi-Lipschitzness for families of Lorentz distances to Cauchy surface foliations. Version arXiv:2606.21997v2 (CC BY 4.0), distributed under the Creative Commons Attribution 4.0 International licence, as recorded in the arXiv licence metadata for this exact version and shown on the abstract page https://arxiv.org/abs/2606.21997v2. Full rights notice: licenses/arxiv-2606.21997-CC-BY-4.0.txt.\par\smallskip

\noindent\textit{Editorial scope.} Counted unit: the Proposition whose source label is \texttt{P:compact splitting implies equi-Lipschitz} in the article (source0.tex lines 467--469, source label \texttt{P:compact splitting implies equi-Lipschitz}), delivered together with the article's own complete original proof of it (source0.tex lines 470--495, one \texttt{proof} environment of 26 lines). No mathematics is written, repaired or re-derived: statement and proof are the article's own text line for line. Required context retained uncounted: Lemma: blowup of initial tangents (lines 357--360). Every definition, display and label of those lines is retained unchanged. Omitted from this package and not redistributed: 1--356, 361--466, 496--519. Nothing inside a retained excerpt is omitted or altered: every retained body is delivered exactly as the article's lines are written, so no omission receipt of any kind accompanies this package. Citations: the retained excerpts carry no citation command at all, so no bibliography entry of the article is transcribed and no reference is redistributed. Reference adaptation: none was needed and none was made; every reference inside the delivered unit resolves inside it, so no unresolved reference or citation is produced. Printed slips kept literal: no printed word, symbol, hypothesis or number of the counted statement or of its proof was altered. Layout and class: the article's own class is \texttt{article} with packages including inputenc, a4wide, amssymb, amsmath, amsthm, euscript, ulem, hyperref, xcolor, appendix, tikz, tikz-cd, todonotes, changebar; the wrapper is the lane's pinned \texttt{amsart} editorial wrapper at 10pt on A4, which restates the theorem-like environments the retained text needs and adds the article's own macro definitions that the retained text needs (\texttt{\textbackslash R}), with extra wrapper packages (bm, bbm, dsfont, xspace, cancel, mathtools). The delivered numbering is the unit's own. Assets: no retained span contains a figure environment, a graphics-inclusion command, a PDF-page inclusion, a table or a TikZ picture; no image file is copied into this package, the package declares no asset dependency and no image file is redistributed with it. Licence: version arXiv:2606.21997v2 of this article is distributed under the Creative Commons Attribution 4.0 International licence, as recorded in the arXiv licence metadata for this exact version and shown on the abstract page https://arxiv.org/abs/2606.21997v2. A full rights notice is delivered at licenses/arxiv-2606.21997-CC-BY-4.0.txt. Characters: every retained excerpt is pure ASCII, so the delivered engine, XeLaTeX with its default Latin Modern text font, reports no missing character.
\begin{Lemma}
\label{Lemma: blowup of initial tangents}
The family $\{\ell(\cdot,\Sigma_t)\}_{t > 0}$ is locally equi-Lipschitz on $M$ if and only if $M$ is covered by open sets $U$ such that for each $q \in U$, the $h$-norms of the initial tangents of $g$-unit speed maximizing timelike geodesics $\gamma_{(t)}$ from $q$ to $\Sigma_t$ satisfy $|\gamma'_{(t)}(0)|_h \leq C$ for some $C > 0$ which depends on $U$ but is independent of $t$, and $t$ is large enough so that $\Sigma_t \subseteq I^+(q)$ for every $q \in U$.
\end{Lemma}

\begin{Proposition}\label{P:compact splitting implies equi-Lipschitz}
    Let $(M,g) = (\R \times S, dt^2 - \tilde g)$ be a product spacetime, where $(S,\tilde g)$ is a compact Riemannian manifold. Then for every Cauchy temporal function $\tau \in C^\infty(\R \times S)$, the families $\{\ell(\cdot,\Sigma_t)\}_{t > 0}$ and $\{\ell(\Sigma_{-t},\cdot)\}_{t > 0}$ are locally equi-Lipschitz.
\end{Proposition}

\begin{proof}
    We show this for $\{\ell(\cdot,\Sigma_t)\}_{t > 0}$, the case of the other family is analogous.

    Given any $s \in \R$, since $\Sigma_s$ is a smooth spacelike acausal Cauchy hypersurface it is easily seen that it must be a global graph over $S$, i.e.,
    \begin{equation}
        \Sigma_s = \{(f_s(x),x) \in \R \times S : x \in S\}
    \end{equation}
    for some smooth function $f_s \in C^\infty(S)$. The fact that $\Sigma_s$ is spacelike translates to the condition $|df_s|_{\tilde g} < 1$, i.e., $f_s$ is $1$-Lipschitz on $S$. Let $D:=\mathrm{diam}(S) < +\infty$ (by compactness), then
    \begin{equation}
        \max_S f_s - \min_S f_s \leq D <+\infty
    \end{equation}
    for every $s \in \R$. We write $M_s:=\max_S f_s$, so that $f_s \geq M_s - D$. Fix a compact set $K \subseteq \R \times S$ and take any $p:=(t_0,x_0) \in K$. For $s$ large enough, $K \subseteq I^-(\Sigma_s)$. Connect $p$ to $\Sigma_s$ by the vertical segment $\alpha:t \mapsto (t + t_0,x_0)$. This segment uniquely meets $\Sigma_s$ at $t = f_s(x_0)$, so that its Lorentzian arclength 
    $ L_g(\alpha)$ is
    \begin{equation}
       \ell(p,\Sigma_s) \geq L_g(\alpha) = f_s(x_0) - t_0 \geq M_s - D - t_0.
    \end{equation}
    Let $\gamma_s(t) = (t,c_s(t))$ be any timelike geodesic realizing $\ell(p,\Sigma_s)$, so that $v_s:=| c'_s|_{\tilde g} < 1$. This geodesic intersects $\Sigma_s$ at some $(f_s(y_s),y_s)$, so that
    \begin{equation}
        \ell(p,\Sigma_s) = L_g(\gamma_s) = (f_s(y_s) - t_0) \sqrt{1 - v_s^2} \leq (M_s - t_0) \sqrt{1 - v_s^2}.
    \end{equation}
    In total,
    \begin{equation}
        M_s - D - t_0 \leq (M_s - t_0) \sqrt{1 - v_s^2}.
    \end{equation}
    Note that $M_s \to +\infty$ as $s \to +\infty$, and moreover $t_0$ remains uniformly bounded for $p = (t_0,x_0) \in K$ since the latter is compact. Thus, dividing by $M_s$ and sending $s \to \infty$, we get $1 \leq \liminf_{s \to \infty} \sqrt{1 - v_s^2}$, so that $\lim_{s \to \infty} v_s^2 \to 0$. This shows that $|\gamma_s'(0)|_h^2 = 1 + v_s^2$ remains uniformly bounded in $s$, where $h:=dt^2 + \tilde g$, thus we conclude the proof upon recalling Lemma \ref{Lemma: blowup of initial tangents}.
\end{proof}

\end{document}
